%% AMS-LaTeX Created with the Wolfram Language : www.wolfram.com

\documentclass[12pt]{article}

\usepackage[letterpaper,margin=1in]{geometry}
\usepackage{amsmath,amssymb,setspace}
\usepackage{graphicx}
\usepackage{bm}
\graphicspath{{./figs/}}
\usepackage{tcolorbox}
\usepackage[framed,numbered,autolinebreaks,useliterate]{mcode}
\usepackage{enumitem}
\usepackage{xcolor}
\usepackage{listings}
\usepackage{hyperref}
\usepackage[cal=boondox]{mathalfa}

% Define colors for lstlisting
\definecolor{codegreen}{rgb}{0,0.6,0}
\definecolor{codegray}{rgb}{0.5,0.5,0.5}
\definecolor{codepurple}{rgb}{0.58,0,0.82}
\definecolor{backcolour}{rgb}{0.95,0.95,0.95}

\lstdefinestyle{mystyle}{
    backgroundcolor=\color{backcolour},   
    commentstyle=\color{codegreen},
    keywordstyle=\color{magenta},
    numberstyle=\tiny\color{codegray},
    stringstyle=\color{codepurple},
    basicstyle=\ttfamily\normalsize,
    breakatwhitespace=false,         
    breaklines=true,                 
    captionpos=b,                    
    keepspaces=true,                 
    numbers=none,  
    columns=flexible,         
    numbersep=5pt,                  
    showspaces=false,                
    showstringspaces=false,
    showtabs=false,                  
    tabsize=2,
    literate={'}{{\textquotesingle}}1 {''}{{\textquotedbl}}2
}
\lstset{style=mystyle}

\hypersetup{
    colorlinks=true,
    linkcolor=blue,
    filecolor=magenta,      
    urlcolor=cyan,
    pdftitle={Your Document Title},
    pdfpagemode=FullScreen,
}

\def\mathbi#1{\textbf{\em #1}}
\newcommand{\fullunderline}{\noindent\makebox[\linewidth]{\rule{\linewidth}{0.1pt}}\\[-2em]}

\begin{document}
\doublespacing	

\title{CE 401/5501 Homework 06\\[2in]
\textbf{Name: \rule{2in}{0.15mm}}}
\author{}
\date{}
\maketitle

\newcommand*{\I}{\imath} % imaginary unit i

\newpage

\section{Multiple Choice / Fill-in-the-blank Questions}

\textbf{Question 1} \\

A perceptron is the simplest artificial neural network, and its mathematical expression can be written as 
\[
y = f\big[u(\bm{x})\big],
\]
where \( u(\bm{x}) = \bm{w}^T \bm{x} + b \). Which of the following statements describes its architecture? 

\begin{itemize}
    \item [$\square$] The functional model describes a computing flow: Input features \(\bm{x}\) \(\rightarrow\) aggregation through \(u(\cdot)\) \(\rightarrow\) activation through \(f(\cdot)\) \(\rightarrow\) outputting a class label or score \(y\).
    \item [$\square$] \(f(\cdot)\) is called an activation function, which is often a simple nonlinear function.
    \item [$\square$] The perceptron model is equivalent to a logistic regression model.
    \item [$\square$] The output \(y\) must be binary.
\end{itemize}

\bigskip

\newpage 

\textbf{Question 2} \\

Given a perceptron model for binary classification, its input is a two-dimensional data point denoted by 
\[
\bm{x} = \begin{bmatrix} x_1 \\ x_2 \end{bmatrix}.
\]
After learning from the data, the resulting learned weights and bias are 
\[
[w_1, w_2, b]^T = [-4, -3, 10]^T.
\]
Answer the following analysis questions:

\begin{itemize}
    \item Write down the decision function \( u(\bm{x}) \).
    \item In the \(x_1\text{-}x_2\) plane, plot the decision boundary \( u(\bm{x}) = 0 \) (hint: write it as \(x_2\) as a function of \(x_1\)).
    \item Given that the activation is a \textit{Heaviside} function, write out the resulting perceptron model.
    \item Given two data points:
    \begin{itemize}
        \item A: \([1, 4]^T\)
        \item B: \([-2, 2]^T\)
    \end{itemize}
    Using the perceptron model, Point A will be labeled as \((\ \ \ \ \ \ )\); and Point B will be labeled as \((\ \ \ \ \ \ )\).
    \item If the activation is a \textit{ReLU} function:
    \begin{itemize}
        \item For Point A: \([1, 4]^T\), the predicted score is \((\ \ \ \ \ \ )\).
        \item For Point B: \([-2, 2]^T\), the predicted score is \((\ \ \ \ \ \ )\).
    \end{itemize}
\end{itemize}

\bigskip

\newpage

\textbf{Question 3} \\

Suppose that a civil engineer is constructing an ML exercise. The input data point is 3-dimensional; namely, each data point is expressed as a \(3\times 1\) vector. There are two class labels, 0 and 1. The engineer has collected 1000 data points and manually labeled them.

The engineer starts with an initial neural network and configures it as follows:
\begin{itemize}
    \item The neural network has two hidden layers. Each hidden layer contains 5 neurons.
    \item The engineer's next task is to determine how many parameters (weights) the NN model has, assuming it is a fully connected multilayer perceptron. Here, we consider bias (intercept) parameters at each layer. Effectively, there is an additional input of 1 at each layer (the input layer and the two hidden layers).
\end{itemize}
Now you are tasked to answer:
\begin{itemize}
    \item How many weight parameters are in the model?
    \item Graph the MLP neural network.
\end{itemize}

\newpage

% \section{Programming Practice and Data Operation}

\end{document}
